\section{OSL-Core: a concrete operator-symbolic language}
\label{sec:osl-core}

The preceding definition describes a family of languages.  This section fixes a
small concrete member of that family, called \OSLCore.  Its purpose is not to
replace neural computation or to settle every metatheoretic question.  It is a
reference intermediate language in which operator declarations, provenance,
proof obligations, memory writes, communication, effects, and action
authorization can be stated in one typed program.

\begin{statusbox}
The release accompanying this manuscript includes a deterministic reference
interpreter for a finite subset, sixty-six source programs, sixty-one automated tests,
and a machine-readable example-suite result.  The interpreter implements the
operational policies described below for finite arrays, provenance-aware
packets, first-class proofs, bounded capabilities, temporal contracts,
transactions, session protocols, neural proposals, memory writes, planning,
effects, and receipts.  It does not implement
the complete higher-order, stochastic, nonlinear, or proof-assistant semantics
proposed elsewhere in the paper.
\end{statusbox}

\subsection{Design goals}

\OSLCore is designed around twelve requirements.
\begin{enumerate}[leftmargin=1.7em]
  \item \textbf{Typed operator composition.}  States and transformations carry
  explicit source and target types.
  \item \textbf{Coordinate tolerance.}  Meaning may be preserved by an
  intertwiner or transported operator even when concrete vectors differ.
  \item \textbf{Provenance conservation.}  Communication preserves source and
  lineage rather than converting foreign content into direct observation.
  \item \textbf{Linear evidence resources.}  Observation and authorization
  tokens cannot be silently forged or duplicated.
  \item \textbf{Effect confinement.}  External action, tool calls, memory
  writes, communication, simulation, and model modification are separate
  effects.
  \item \textbf{Safe rejection.}  A program that lacks a proof or capability is
  rejected or quarantined rather than implicitly coerced into validity.
  \item \textbf{Canonical receipts.}  Every state-changing instruction emits a
  canonical record whose digest binds the declared transition.
  \item \textbf{Hybrid execution.}  Learned latent modules may propose states
  and operators, while the typed verifier controls which proposals may become
  memory or action.
  \item \textbf{First-class proofs.}  Access, origin, clearance, endorsement,
  authorization, ownership, and use may be materialized as digest-bound proof
  terms rather than remaining implicit runtime flags.
  \item \textbf{Temporal contracts.}  Safety invariants and bounded response
  obligations are monitored over the execution trace.
  \item \textbf{Transactional repair.}  Memory and model updates may be
  checkpointed, committed, or rolled back atomically in the finite runtime.
  \item \textbf{Protocol-typed sharing.}  Finite-state sessions constrain the
  order and participants of shared-workspace communication.
\end{enumerate}

\subsection{Lexical conventions}

Identifiers begin with a letter or underscore and may contain letters, digits,
underscores, and hyphens.  Comments begin with \texttt{\#}.  Strings are
UTF-8 double-quoted literals.  Numeric arrays use JSON-compatible bracket
notation.  Every source statement occupies one logical line in the reference
subset, although a production parser would permit ordinary block layout.

The language is case-sensitive.  Reserved words are shown in
\cref{lst:osl-keywords}.  Canonical serialization normalizes field order,
numeric precision, set order, and whitespace before hashing; textual spelling
is therefore not itself a trusted provenance source.

\begin{lstlisting}[style=osl,caption={Reserved words in the finite reference subset.},label={lst:osl-keywords}]
module lineage extends agent grant effect
space dim vector in values operator kind matrix
apply transport check_projector check_intertwiner between and tolerance
observe infer simulate inject communicate from to as
bind quarantine endorse threshold authorize for write_memory
plan using expected execute feedback observed
proof mode persistent affine linear requires attest neural_propose admit reject_candidate
contract after within check_contracts begin_transaction commit_transaction rollback_transaction
session participants initial session_transition session_terminal verify_session start_session
secret metadata receipt_check verify_receipts
assert packet provenance accessible_to quarantined_by
memory contains action executed vector approx parent show
\end{lstlisting}

\subsection{Concrete type grammar}

The abstract type universe is refined to the following concrete family:
\[
\begin{aligned}
\tau ::= {}&
\texttt{Bool}
\mid \texttt{Nat}
\mid \texttt{Real}
\mid \texttt{String}
\mid \texttt{Agent}
\mid \texttt{Lineage}\\
&\mid \texttt{State}[n]
\mid \texttt{Tensor}[d_1,\ldots,d_k]
\mid \texttt{Operator}[\tau\Rightarrow\sigma]\\
&\mid \texttt{Packet}[\kappa,\tau]
\mid \texttt{Workspace}[a,\tau]
\mid \texttt{Memory}[a,\kappa,\tau]\\
&\mid \texttt{Proof}[\varphi]
\mid \texttt{Action}[\mathcal E]
\mid \texttt{Receipt}[\varphi]
\mid \texttt{Distribution}[\tau].
\end{aligned}
\]
Here \(n,d_i\in\mathbb N_{>0}\), \(a\) is an agent, \(\kappa\) is a
provenance class, and \(\mathcal E\) is an effect row.

\begin{longtable}{p{0.26\textwidth}p{0.30\textwidth}p{0.34\textwidth}}
\toprule
Type & Static role & Runtime interpretation \\
\midrule
\endfirsthead
\toprule
Type & Static role & Runtime interpretation \\
\midrule
\endhead
\texttt{State[n]} & finite causal or latent state & vector in \(\mathbb R^n\) in the reference backend \\
\texttt{Tensor[d1,...,dk]} & higher-order structured state & multidimensional array or tensor-network handle \\
\texttt{Operator[t=>s]} & typed transformation & matrix, nonlinear map, symbolic operator, or verified external primitive \\
\texttt{Packet[k,t]} & content with provenance class \(k\) & canonical content, source, lineage, confidence, authorization, digest \\
\texttt{Workspace[a,t]} & content accessible to agent \(a\) & packet admitted to the private or shared workspace \\
\texttt{Memory[a,k,t]} & autobiographical record of class \(k\) & immutable record linked to a source packet digest \\
\texttt{Proof[p]} & evidence for proposition \(p\) & proof term, capability, or verifier receipt \\
\texttt{Action[E]} & action whose effects are in \(E\) & planned transition awaiting capability-checked execution \\
\texttt{Receipt[p]} & canonical transition witness & integrity-bound record supporting proposition \(p\) \\
\bottomrule
\end{longtable}

\subsection{Provenance classes}

The concrete provenance class is
\[
\kappa\in
\left\{
\begin{array}{l}
\mathsf{observed},\mathsf{inferred},\mathsf{recalled},
\mathsf{simulated},\mathsf{communicated},\mathsf{injected},\\
\mathsf{foreign},\mathsf{reconstructed},\mathsf{conflicted},
\mathsf{unresolved}
\end{array}
\right\}.
\]
These tags are not mere comments.  They constrain the constructors available
for memory and planning.  The same decoded content may inhabit distinct packet
types:
\[
\texttt{Packet[observed,RouteStatus]}
\quad\text{and}\quad
\texttt{Packet[injected,RouteStatus]}.
\]
The types share a content carrier but have different admissible eliminations.

\subsection{Resource qualifiers}

A production language should distinguish at least three resource modes:
\[
q::=\mathsf{persistent}\mid\mathsf{affine}\mid\mathsf{linear}.
\]
Persistent values may be copied.  Affine values may be used at most once.
Linear values must be consumed exactly once unless an explicit escrow or
cancellation proof is supplied.  Direct observation tokens, one-shot
capabilities, and cryptographic nonces are linear.  Content strings and public
operator declarations are normally persistent.

The full static judgment is
\[
\Gamma;\Delta;\Pi;\Omega
\vdash e:\tau
\dashv\Delta'
\;!\;\mathcal E,
\]
where:
\begin{itemize}[leftmargin=1.6em]
  \item \(\Gamma\) contains persistent variables and declarations;
  \item \(\Delta\) contains affine or linear evidence resources;
  \item \(\Pi\) contains proved propositions and trusted policy assumptions;
  \item \(\Omega\) contains effect capabilities;
  \item \(\Delta'\) records unconsumed resources;
  \item \(\mathcal E\) is the emitted effect row.
\end{itemize}
The finite Python interpreter enforces the most important resource distinction
operationally: a communicated packet contains an origin record but no new
observation token for the recipient.

\subsection{Effect rows and capabilities}

Effect rows are finite sets drawn from
\[
\Effects=
\{
\mathsf{read},
\mathsf{writeMemory},
\mathsf{communicate},
\mathsf{toolCall},
\mathsf{externalAct},
\mathsf{simulate},
\mathsf{modifyModel}
\}.
\]
An effect may be inferred statically, authorized symbolically, and separately
granted as a runtime capability.  These are different gates:
\[
\operatorname{WellTyped}
\not\Rightarrow
\operatorname{Authorized}
\not\Rightarrow
\operatorname{CapabilityPresent}.
\]
This three-stage separation is deliberate.  A plan may be mathematically
well-formed but still lack permission or execution capability.

\subsection{Surface grammar}

\Cref{lst:osl-ebnf} gives the declarative grammar for the reference subset.  It
is intentionally compact enough to inspect manually.  A larger implementation
would add parametric polymorphism, modules, recursive data, stochastic kernels,
pattern matching, asynchronous channels, and verified foreign-function
interfaces.

\lstinputlisting[style=osl,caption={EBNF for OSL-Core 0.3.},label={lst:osl-ebnf}]{reference/osl_core.ebnf}

\subsection{Declaration forms}

A source file begins by declaring a module, lineages, agents, spaces, and
operators.  For example:
\begin{lstlisting}[style=osl]
module mapping_agent
lineage L_root
lineage L_mobile extends L_root
agent Mapper lineage L_mobile
grant Mapper effect externalAct
space Hidden dim 3
operator Pself : Hidden -> Hidden kind projector
  matrix [[1,0,0],[0,1,0],[0,0,0]]
\end{lstlisting}
A lineage declaration creates an explicit parent relation.  It does not by
itself prove identity preservation; it only creates the object on which later
transformation and receipt proofs may operate.

\subsection{Operator forms}

The concrete subset recognizes four operator roles:
\[
\mathsf{linear},\quad
\mathsf{projector},\quad
\mathsf{transport},\quad
\mathsf{external}.
\]
The interpreter implements finite linear and projector matrices.  A projector
declaration is accepted only when
\[
\|P^2-P\|\leq\tau_P.
\]
Transport of an endomorphism is compiled to
\[
P'=TPT^{-1}.
\]
The command
\begin{lstlisting}[style=osl]
check_intertwiner T between A and B tolerance 1e-9
\end{lstlisting}
checks
\[
\|TA-BT\|\leq 10^{-9}.
\]
A future nonlinear backend would replace matrices by typed callable objects and
would check an approximate conjugacy or bisimulation obligation.

\subsection{Packet constructors}

The four primitive event constructors have distinct semantics:
\begin{align*}
\operatorname{observe}_a(c,q)
&:\texttt{Packet[observed,Content]},\\
\operatorname{infer}_a(c\mid p,q)
&:\texttt{Packet[inferred,Content]},\\
\operatorname{simulate}_a(c,q)
&:\texttt{Packet[simulated,Content]},\\
\operatorname{inject}_{\mathrm{ext}\to a}(c,q)
&:\texttt{Packet[injected,Content]}.
\end{align*}
Only the observation constructor creates the linear token
\(\Obs_a(c)\).  Simulation creates a self-generated source record but no
observation token.  Injection creates workspace access and immediate
quarantine, but no self-binding or ownership authorization.

\subsection{Communication}

Communication copies content but transports provenance:
\[
\frac{
  p:\texttt{Packet}[\kappa,\tau]
  \qquad
  \Access_A(p)
}{
  \operatorname{send}_{A\to B}(p):
  \texttt{Packet}[\mathsf{communicated},\tau]
}.
\]
The result contains
\[
\Origin_A(p)\otimes\Access_B(p),
\]
not \(\Obs_B(p)\).  Repeated forwarding retains the original source and a
parent-digest chain; it does not rewrite the content as a direct observation by
each intermediate agent.

\subsection{Binding, endorsement, and authorization}

The language separates three judgments that are often conflated:
\begin{enumerate}[leftmargin=1.7em]
  \item \texttt{bind} associates accessible content with an agent's active
  planning context or lineage;
  \item \texttt{endorse} checks whether confidence exceeds a declared threshold
  and whether quarantine is absent;
  \item \texttt{authorize} grants a named use such as planning, communication,
  or a particular effect.
\end{enumerate}
Thus:
\[
\operatorname{Bound}(p,a)
\not\Rightarrow
\operatorname{Endorsed}(p,a)
\not\Rightarrow
\operatorname{Authorized}(p,a,\chi).
\]
The distinctions permit the system to maintain an accessible but disputed
claim, a believed but non-actionable claim, or an actionable claim whose effect
is still blocked by the capability layer.

\subsection{Memory writes}

A memory instruction has the form
\begin{lstlisting}[style=osl]
write_memory event_at_A by A as communicated
\end{lstlisting}
and compiles to a proof obligation:
\[
\operatorname{Accessible}_A(p)
\wedge
\operatorname{Bound}_A(p)
\wedge
\operatorname{Compatible}(\rho(p),\kappa)
\wedge
\neg\operatorname{Quarantined}_A(p).
\]
For \(\kappa=\mathsf{observed}\), it additionally requires a nonforgeable
observation token matching the active agent and lineage.  A communicated packet
may be stored as communicated knowledge, but not as the receiver's direct
observation.

\subsection{Planning and execution}

Planning creates a typed action but does not execute it:
\begin{lstlisting}[style=osl]
plan alternate_route by A using route_blocked effect externalAct expected 1.0
\end{lstlisting}
The source packet must be accessible, bound, endorsed, authorized for planning
or the named effect, and unquarantined.  Execution is a second transition:
\begin{lstlisting}[style=osl]
execute alternate_route by A
\end{lstlisting}
which checks that the agent possesses the runtime capability
\(\mathsf{externalAct}\).  This separation permits dry-run planning,
third-party verification, human approval, or sandbox evaluation before any
external effect occurs.

\subsection{Feedback}

The command
\begin{lstlisting}[style=osl]
feedback alternate_route by A observed 1.2
\end{lstlisting}
links the observed consequence to the action receipt.  If an expected value was
recorded, the runtime also emits the prediction error
\[
e=r_{\mathrm{observed}}-r_{\mathrm{expected}}.
\]
The finite core records the error but does not prescribe one universal learning
rule.  A recurrent agent may use it to update confidence, request more evidence,
retrain an operator, or reopen a quarantined hypothesis.

\subsection{Assertions and executable specification}

Packaged examples contain executable assertions such as:
\begin{lstlisting}[style=osl]
assert packet route_at_A provenance communicated
assert memory A contains route_at_A as observed false
assert action alternate_route executed true
assert vector self_x approx [2.0,-1.0,0.0] tolerance 1e-12
\end{lstlisting}
These assertions are part of the example harness, not the trusted runtime
language.  They state expected postconditions and provide a compact executable
specification of the intended semantics.

\subsection{Canonical packet representation}

A runtime packet is the canonical tuple
\[
p=
(c,a_{\mathrm{src}},L_{\mathrm{src}},\kappa,q,
A,B,Q,K,\mathcal U,\mathcal O,d_{\mathrm{parent}},d_p),
\]
where:
\begin{itemize}[leftmargin=1.6em]
  \item \(A\) is the set of agents with access;
  \item \(B\) is the set of agents to which the packet is currently bound;
  \item \(Q\) is the quarantine set;
  \item \(K\) is the endorsement set;
  \item \(\mathcal U\) is the authorization map;
  \item \(\mathcal O\) is the set of direct-observation tokens;
  \item \(d_{\mathrm{parent}}\) links derived or communicated content to its
  source packet;
  \item \(d_p\) is the canonical packet digest.
\end{itemize}
The digest establishes integrity of the canonical representation.  It does not
establish authorship, truth, or semantic authorization.
